% Options for packages loaded elsewhere
\PassOptionsToPackage{unicode}{hyperref}
\PassOptionsToPackage{hyphens}{url}
\PassOptionsToPackage{dvipsnames,svgnames,x11names}{xcolor}
\documentclass[
  10pt,
]{article}
\usepackage{xcolor}
\usepackage[margin=2.4cm]{geometry}
\usepackage{amsmath,amssymb}
\setcounter{secnumdepth}{-\maxdimen} % remove section numbering
\usepackage{iftex}
\ifPDFTeX
  \usepackage[T1]{fontenc}
  \usepackage[utf8]{inputenc}
  \usepackage{textcomp} % provide euro and other symbols
\else % if luatex or xetex
  \usepackage{unicode-math} % this also loads fontspec
  \defaultfontfeatures{Scale=MatchLowercase}
  \defaultfontfeatures[\rmfamily]{Ligatures=TeX,Scale=1}
\fi
\usepackage{lmodern}
\ifPDFTeX\else
  % xetex/luatex font selection
  \setmainfont[]{Libertinus Serif}
  \setmonofont[Scale=0.85]{Latin Modern Mono}
  \setmathfont[]{Latin Modern Math}
\fi
% Use upquote if available, for straight quotes in verbatim environments
\IfFileExists{upquote.sty}{\usepackage{upquote}}{}
\IfFileExists{microtype.sty}{% use microtype if available
  \usepackage[]{microtype}
  \UseMicrotypeSet[protrusion]{basicmath} % disable protrusion for tt fonts
}{}
\makeatletter
\@ifundefined{KOMAClassName}{% if non-KOMA class
  \IfFileExists{parskip.sty}{%
    \usepackage{parskip}
  }{% else
    \setlength{\parindent}{0pt}
    \setlength{\parskip}{6pt plus 2pt minus 1pt}}
}{% if KOMA class
  \KOMAoptions{parskip=half}}
\makeatother
\usepackage{longtable,booktabs,array}
\newcounter{none} % for unnumbered tables
\usepackage{calc} % for calculating minipage widths
% Correct order of tables after \paragraph or \subparagraph
\usepackage{etoolbox}
\makeatletter
\patchcmd\longtable{\par}{\if@noskipsec\mbox{}\fi\par}{}{}
\makeatother
% Allow footnotes in longtable head/foot
\IfFileExists{footnotehyper.sty}{\usepackage{footnotehyper}}{\usepackage{footnote}}
\makesavenoteenv{longtable}
\setlength{\emergencystretch}{3em} % prevent overfull lines
\providecommand{\tightlist}{%
  \setlength{\itemsep}{0pt}\setlength{\parskip}{0pt}}

\providecommand{\E}{\mathbb{E}}
\providecommand{\N}{\mathbb{N}}
\providecommand{\Z}{\mathbb{Z}}
\providecommand{\R}{\mathbb{R}}
\providecommand{\eps}{\varepsilon}
\providecommand{\ceil}[1]{\left\lceil #1\right\rceil}
\providecommand{\floor}[1]{\left\lfloor #1\right\rfloor}
\AtBeginDocument{\let\setminus\smallsetminus}
\usepackage[most]{tcolorbox}
\newtcolorbox{warning}{colback=red!5,colframe=red!65!black,boxrule=0.8pt,arc=2pt,left=6pt,right=6pt,top=5pt,bottom=5pt}
\usepackage{etoolbox}
\AtBeginEnvironment{longtable}{\small}
\setlength{\emergencystretch}{3em}
\usepackage{fancyhdr}
\pagestyle{fancy}
\fancyhf{}
\fancyhead[L]{\small\bfseries UNREFEREED GPT-6 ASTRA OUTPUT}
\fancyhead[R]{\small\thepage}
\renewcommand{\headrulewidth}{0.4pt}
\usepackage{bookmark}
\IfFileExists{xurl.sty}{\usepackage{xurl}}{} % add URL line breaks if available
\urlstyle{same}
\hypersetup{
  pdftitle={Mixing Circular Colourings},
  pdfauthor={Writeup: gpt-6-astra (single model pass)},
  colorlinks=true,
  linkcolor={blue!50!black},
  filecolor={Maroon},
  citecolor={Blue},
  urlcolor={blue!50!black},
  pdfcreator={LaTeX via pandoc}}

\title{Mixing Circular Colourings}
\usepackage{etoolbox}
\makeatletter
\providecommand{\subtitle}[1]{% add subtitle to \maketitle
  \apptocmd{\@title}{\par {\large #1 \par}}{}{}
}
\makeatother
\subtitle{Unrefereed candidate proof by GPT-6 Astra}
\author{Writeup: \texttt{gpt-6-astra} (single model pass)}
\date{Catalog id \texttt{mixing\_circular\_colourings\_0} --- generated
2026-09-16}

\begin{document}
\maketitle

\begin{warning}

\textbf{UNREFEREED MODEL OUTPUT.} This document was selected solely
because GPT-6 Astra labelled its own result
\texttt{would\_publish:\ true} and returned \texttt{proved} or
\texttt{disproved}. It has not passed the adversarial LLM referee used
for the earlier Sol campaign, has not been checked by a human
mathematician, and has not been checked for novelty. Treat every
mathematical and bibliographic claim below as unverified.

\end{warning}

\section{Summary and provenance}\label{summary-and-provenance}

{\def\LTcaptype{none} % do not increment counter
\begin{longtable}[]{@{}
  >{\raggedright\arraybackslash}p{(\linewidth - 2\tabcolsep) * \real{0.5000}}
  >{\raggedright\arraybackslash}p{(\linewidth - 2\tabcolsep) * \real{0.5000}}@{}}
\toprule\noalign{}
\endhead
\bottomrule\noalign{}
\endlastfoot
Catalog id & \texttt{mixing\_circular\_colourings\_0} \\
Catalog entry &
\href{https://graph-theory-ai.github.io/graph-conjectures/op/mixing_circular_colourings_0/}{Mixing
Circular Colourings} \\
Source corpus & opg \\
Campaign leg & OpenProblemGarden first pass (\texttt{attacks\_opg}) \\
Source paper / entry &
http://www.openproblemgarden.org/op/mixing\_circular\_colourings\_0 \\
Model verdict & \textbf{proved} (confidence: high) \\
Model's one-line claim & For an n-vertex graph with an edge, the
circular mixing threshold is rational, with reduced numerator at most
n+1. \\
Model & \texttt{gpt-6-astra}, reasoning effort \texttt{max},
\texttt{mode=pro}, flex service tier \\
Original artifact &
\texttt{attacks\_opg/mixing\_circular\_colourings\_0/output.md} \\
Independent review & \textbf{None} \\
Model's caveats & Finite loopless graphs; edgeless graphs require the
usual lower-end normalization. Integrality and attainment are not
established. \\
\end{longtable}
}

\section{Problem statement}\label{problem-statement}

Question Is \(\mathfrak{M}_c(G)\) always rational?

\subsection{Catalog context}\label{catalog-context}

Given a proper \(k\) -colouring \(f\) of a graph \(G\) , consider the
following `mixing process:'

\begin{itemize}
\item
  choose a vertex \(v\in V(G)\) ;
\item
  change the colour of \(v\) (if possible) to yield a different \(k\)
  -colouring \(f'\) of \(G\) . A natural problem arises: Can every \(k\)
  -colouring of \(G\) be generated from \(f\) by repeatedly applying
  this process? If so, we say that \(G\) is \(k\) -mixing . The problem
  of determining if a graph is \(k\) -mixing and several related
  problems have been studied in a series of recent papers {[}1,2,4-6{]}.
  The authors of {[}4{]} provide examples which show that a graph can be
  \(k\) -mixing but not \(k'\) -mixing for integers \(k' > k\) . For
  example, given \(m\geq3\) consider the bipartite graph \(L_m\) which
  is obtained by deleting a perfect matching from \(K_{m,m}\) . It is an
  easy exercise to show that for \(L_m\) is \(k\) -mixing if and only if
  \(k\geq3\) and \(k\neq m\) . This example motivates the following
  definition. Definition Define the mixing threshold of \(G\) to be
  \[\mathfrak{M}(G) := \min\{\ell\in\mathbb{N}: G\text{ is }k\text{-mixing whenever } k\geq\ell\}.\]
  An analogous definition can be made for circular colouring. Recall, a
  \((k,q)\) -colouring of a graph is a mapping
  \(f:V(G)\to \{0,1,\dots,k-1\}\) such that if \(uv\in E(G)\) , then
  \(q\leq |f(u)-f(v)|\leq k-q\) . As with ordinary colourings, we say
  that a graph \(G\) is \((k,q)\) -mixing if all \((k,q)\) -colourings
  of \(G\) can be generated from a single \((k,q)\) -colouring \(f\) by
  recolouring one vertex at a time. Definition Define the circular
  mixing threshold of \(G\) to be
  \[\mathfrak{M}_c(G) := \inf\{\ell\in\mathbb{Q}: G\text{ is }(k,q)\text{-mixing whenever } k/q \geq\ell\}.\]
  Several bounds on the circular mixing threshold are obtained in
  {[}3{]}, including the following which relates the circular mixing
  threshold to the mixing threshold. Theorem (Brewster and Noel
  \([3)</b>\) {]} For every graph \(G\) ,
  \[\mathfrak{M}_c(G)\leq\max\left\{\frac{|V(G)|+1}{2}, \mathfrak{M}(G)\right\}.\]
  As a corollary, we have the following: if
  \(|V(G)|\leq 2\mathfrak{M}(G)-1\) , then
  \(\mathfrak{M}_c(G)\leq \mathfrak{M}(G)\) . However, examples in
  {[}3{]} show that the ratio \(\mathfrak{M}_c/\mathfrak{M}\) can be
  arbitrarily large in general. Regarding the problem of determining if
  \(\mathfrak{M}_c\) is rational, it is worth mentioning that there are
  no known examples of graphs \(G\) for which \(\mathfrak{M}_c(G)\) is
  not an integer. Other problems are also given in {[}3{]}. One can
  check that \(\mathfrak{M}_c(K_2) = 2\) and \(\mathfrak{M}(K_2) = 3\) .
  However, the only graphs which are known to satisfy
  \(\mathfrak{M}_c < \mathfrak{M}\) are in some sense related to \(K_2\)
  , eg. trees and complete bipartite graphs. Question Is there a
  non-bipartite graph \(G\) such that
  \(\mathfrak{M}_c(G) < \mathfrak{M}(G)\) ? Using an example from
  {[}4{]}, it is shown in {[}3{]} that if \(m\geq2\) is an integer, then
  there is a graph \(G\) such that \(\mathfrak{M}_c(G) = \chi(G) = m\)
  if and only if \(m\neq 3\) . A natural question to ask is whether a
  similar result holds for the circular chromatic number. Again, certain
  bipartite graphs are an exception. Question Is there a non-bipartite
  graph \(G\) such that \(\mathfrak{M}_c(G) =\chi_c(G)\) ? Also, the
  example of \(K_2\) shows that the circular mixing threshold is, in
  general, not attained. However, the following problem is open.
  Question Is the circular mixing threshold always attained for
  non-bipartite graphs? For more precise versions of the last three
  questions, see {[}3{]}.
\end{itemize}

\section{Astra writeup}\label{astra-writeup}

\emph{The text below is the model output verbatim, apart from moving its
machine-readable verdict block into the summary above and shifting
Markdown heading levels for this document. Mechanical TeX defects and
equation tags were normalized where needed for compilation.}

\subsection{A finite set of possible circular mixing
thresholds}\label{a-finite-set-of-possible-circular-mixing-thresholds}

Write \(\mathcal R_{k,q}(G)\) for the graph whose vertices are the
\((k,q)\)-colourings of \(G\), with two colourings adjacent when they
differ at exactly one vertex. Thus mixing means that this graph is
nonempty and connected.

For \(N\ge 2\), define \[
\mathcal E_N
=
\left\{\frac ab:\ a,b\in\mathbb Z_{>0},\quad 2b\le a\le N\right\}.
\]

I prove the following stronger statement.

\textbf{Theorem.} Let \(G\) be a finite loopless graph with \(n\)
vertices and at least one edge. Then \[
\boxed{\mathfrak M_c(G)\in\mathcal E_{n+1}.}
\] In particular, \(\mathfrak M_c(G)\) is rational, and its reduced
numerator is at most \(n+1\).

The main point is that, outside a finite rational set, colourings with a
fixed orientation of the edges form a connected reconfiguration class.
Connections between these classes can then be tested using systems of
difference constraints on at most \(n+1\) variables.

\subsubsection{1. Connectivity of integer difference-constraint
regions}\label{connectivity-of-integer-difference-constraint-regions}

We first need an elementary lemma.

\textbf{Lemma 1.} Let \[
P=\left\{x\in\mathbb Z^s:
0\le x_i\le K,\quad
x_j-x_i\le c_{ij}\text{ for }(i,j)\in A
\right\},
\] where all \(c_{ij}\) and \(K\) are integers. Regard \(A\) as a
directed graph, with arc \(i\to j\) having weight \(c_{ij}\).

Suppose \(P\ne\varnothing\), and no directed simple cycle in \(A\) has
total weight zero. Then the graph on \(P\) allowing a change of \(1\) or
\(-1\) in one coordinate at a time is connected.

\textbf{Proof.} First, \(P\) is closed under coordinatewise minimum.
Indeed, for \(x,y\in P\), put \(z_i=\min(x_i,y_i)\). If \(z_i=x_i\),
then \[
z_j-z_i\le x_j-x_i\le c_{ij};
\] the other case follows using \(y\).

It suffices to connect \(x\) to \(z=\min(x,y)\) by successive coordinate
decreases. At a current point \(x\ge z\), let \[
D=\{i:x_i>z_i\}.
\] A coordinate \(i\in D\) can be decreased by one unless there is a
tight outgoing constraint \[
x_j-x_i=c_{ij}.
\] The box constraints cannot prevent this decrease because
\(x_i>z_i\ge0\).

Moreover, every such blocking coordinate \(j\) belongs to \(D\).
Otherwise \(x_j=z_j\), so \[
z_j-z_i\ge x_j-(x_i-1)=c_{ij}+1,
\] contrary to \(z\in P\).

If every coordinate in \(D\) were blocked, choosing a blocking arc at
each vertex of \(D\) would produce a directed cycle of tight
constraints. Summing around it would give total weight zero, a
contradiction.

Consequently some coordinate can be decreased. Integrality ensures that
every nontight constraint has sufficient slack for a unit decrease.
Repeating reaches \(z\), since \(\sum_i(x_i-z_i)\) strictly decreases.
Applying the same argument to \(y\) completes the proof. \(\square\)

\subsubsection{2. Partitioning colourings by edge
orientations}\label{partitioning-colourings-by-edge-orientations}

Fix \(k,q\in\mathbb Z_{>0}\) with \(k\ge2q\). Every \((k,q)\)-colouring
\(f\) induces an orientation \(D_f\) of \(G\): orient \(uv\) from \(u\)
to \(v\) when \(f(u)<f(v)\).

For an orientation \(D\), let \(X_D(k,q)\) be the colourings inducing
\(D\). These sets partition the colourings of \(G\). If \(u\to v\) is an
edge of \(D\), the corresponding constraints are \[
x_u-x_v\le -q,
\qquad
x_v-x_u\le k-q,
\] together with \(0\le x_v\le k-1\).

The variable-to-variable constraint graph therefore has arc weights only
of the forms \[
-q,\qquad k-q.
\] A directed simple cycle of length \(b\), containing \(a\) arcs of
weight \(k-q\), has total weight \[
a(k-q)+(b-a)(-q)=ak-bq.
\] If this is zero, then \(a\ge1\) and \[
\frac{k}{q}=\frac ba.
\] Since \(b\le n\) and \(k/q\ge2\), this ratio belongs to
\(\mathcal E_n\).

Lemma 1 consequently gives:

\textbf{Lemma 2.} If \(k/q\notin\mathcal E_n\), every nonempty
\(X_D(k,q)\) is connected by single-vertex recolourings. In fact,
changes of one unit in the numerical colour suffice within \(X_D(k,q)\).

The exclusion of \(\mathcal E_n\) is important: zero-weight cycles can
force several coordinates to move together, preventing single-coordinate
connectivity.

\subsubsection{3. Feasibility has only finitely many rational
breakpoints}\label{feasibility-has-only-finitely-many-rational-breakpoints}

We next determine how the nonemptiness of \(X_D(k,q)\) can depend on the
parameters.

\textbf{Lemma 3.} Fix an oriented graph \(D\) on \(s\) vertices. The
nonemptiness of \(X_D(k,q)\):

\begin{enumerate}
\def\labelenumi{\arabic{enumi}.}
\tightlist
\item
  depends only on the ratio \(k/q\), not on its integer representation;
  and
\item
  is constant on each open interval in \[
  (2,\infty)\setminus\mathcal E_s.
  \]
\end{enumerate}

\textbf{Proof.} Use the difference-constraint graph above, and add a
reference vertex \(*\), with arcs \[
*\longrightarrow i \text{ of weight }k-1,
\qquad
i\longrightarrow * \text{ of weight }0.
\] Setting \(x_*=0\) imposes \(0\le x_i\le k-1\).

This integer difference-constraint system is feasible if and only if the
augmented graph has no negative directed cycle. For completeness,
sufficiency follows by taking shortest-path distances from \(*\). All
vertices are reachable; without a negative cycle, these distances are
finite integers and satisfy every constraint. The arcs incident with
\(*\) ensure that the distances lie between \(0\) and \(k-1\).

It is enough to check simple directed cycles. There are two kinds.

\begin{itemize}
\item
  A cycle avoiding \(*\) has weight \[
  ak-bq,\qquad 1\le b\le s,\quad 0\le a\le b.
  \] If \(a=0\), it is always negative. Otherwise, nonnegativity is
  equivalent to \[
  \frac{k}{q}\ge\frac ba.
  \]
\item
  A cycle containing \(*\) uses exactly one arc of weight \(k-1\) and
  one arc of weight \(0\). Its weight is \[
  ak-bq-1,\qquad 0\le b\le s-1,\quad 1\le a\le b+1.
  \] Since \(ak-bq\) is an integer, \[
  ak-bq-1\ge0
  \quad\Longleftrightarrow\quad
  ak-bq>0
  \quad\Longleftrightarrow\quad
  \frac{k}{q}>\frac ba.
  \]
\end{itemize}

Thus feasibility is determined by finitely many weak or strict
comparisons of \(k/q\) with fixed rational numbers. Every comparison
point lying in \([2,\infty)\) belongs to \(\mathcal E_s\). This proves
both assertions. \(\square\)

The integrality step \[
ak-bq-1\ge0\iff ak-bq>0
\] is what removes any apparent dependence on the scale of the pair
\((k,q)\).

\subsubsection{4. Testing transitions between orientation
classes}\label{testing-transitions-between-orientation-classes}

For fixed \(k,q\), construct a finite graph \(Q_{k,q}(G)\) as follows.

\begin{itemize}
\tightlist
\item
  Its vertices are the orientations \(D\) for which
  \(X_D(k,q)\ne\varnothing\).
\item
  Two distinct orientations \(D,D'\) are adjacent if some colouring in
  \(X_D(k,q)\) and some colouring in \(X_{D'}(k,q)\) differ at exactly
  one vertex.
\end{itemize}

We can test each possible adjacency using another oriented graph, now
with \(n+1\) vertices.

Suppose the changed vertex is \(v\). Necessarily \(D,D'\) agree on every
edge not incident with \(v\). Replace \(v\) by two nonadjacent twins
\(v_0,v_1\), each having the original neighbourhood of \(v\). Orient:

\begin{itemize}
\tightlist
\item
  the edges at \(v_0\) as in \(D\);
\item
  the edges at \(v_1\) as in \(D'\);
\item
  all other edges in their common direction.
\end{itemize}

A colouring respecting this orientation specifies the unchanged colours
on \(V(G)\setminus\{v\}\), the old colour at \(v_0\), and the new colour
at \(v_1\). It therefore gives exactly the desired transition.

Conversely, every such transition yields a colouring of this oriented
twin graph. Because \(D\ne D'\), at least one incident edge is
oppositely oriented, so the two twin colours are automatically
different.

Thus adjacency in \(Q_{k,q}(G)\) is determined by finitely many
feasibility tests of the type in Lemma 3, on \(n+1\) vertices.

When \(k/q\notin\mathcal E_n\), Lemma 2 says that all nonempty
orientation classes are connected. Hence \[
\mathcal R_{k,q}(G)\text{ is nonempty and connected}
\quad\Longleftrightarrow\quad
Q_{k,q}(G)\text{ is nonempty and connected}.
\] Indeed, a connected quotient can be lifted by moving within each
class to the endpoint of the next transition; the converse follows by
projecting a reconfiguration path.

We have proved the central stability statement.

\textbf{Proposition 4.} Let \(I\) be a component interval of \[
(2,\infty)\setminus\mathcal E_{n+1}.
\] Then either \(G\) is \((k,q)\)-mixing for every pair of positive
integers with \(k/q\in I\), or it is \((k,q)\)-mixing for none of those
pairs.

\textbf{Proof.} Throughout \(I\), Lemma 3 makes every vertex-feasibility
and transition-feasibility test defining \(Q_{k,q}(G)\) constant. Thus
the finite graph \(Q_{k,q}(G)\) is unchanged.

Also \(I\cap\mathcal E_n=\varnothing\), so all its nonempty orientation
classes are connected by Lemma 2. The quotient criterion now proves the
assertion. \(\square\)

No assertion of scale invariance at the exceptional ratios is needed.

\subsubsection{5. A self-contained upper
bound}\label{a-self-contained-upper-bound}

We also need to know that all sufficiently large ratios are mixing. The
same machinery gives \[
G\text{ is }(k,q)\text{-mixing whenever }k/q\ge n+1.
\]

Indeed, the palette \[
C=\{0,q,2q,\ldots,nq\}
\] has \(n+1\) colours, and every two distinct palette colours satisfy
the \((k,q)\)-constraint when \(k\ge(n+1)q\).

For any colouring \(f\), its induced orientation \(D_f\) is acyclic,
since colours strictly increase along directed edges. A topological
ordering therefore gives an injective palette colouring belonging to
\(X_{D_f}(k,q)\). Because \(k/q\ge n+1\) lies outside \(\mathcal E_n\),
Lemma 2 connects \(f\) to this palette colouring.

Finally, all injective maps from \(V(G)\) into \(C\) are connected by
single-vertex moves that preserve injectivity. To reach a target
injective map, fix vertices one at a time. If a desired colour is
occupied by another, not-yet-fixed vertex, first move that vertex to an
unused palette colour. An unused colour always exists because there are
\(n+1\) palette colours and only \(n\) vertices.

All these intermediate maps are proper circular colourings. This proves
the upper bound, including nonemptiness.

\subsubsection{6. Extracting the
threshold}\label{extracting-the-threshold}

Define the set of bad ratios by \[
B_G=
\left\{
r\in\mathbb Q_{\ge2}:
\begin{array}{l}
\text{there exist }k,q\in\mathbb Z_{>0}\text{ with }k/q=r\\
\text{such that }G\text{ is not }(k,q)\text{-mixing}
\end{array}
\right\}.
\] The existential quantifier accommodates possible dependence on the
integer representation at exceptional ratios.

Because \(G\) has an edge, it is not \((2,1)\)-mixing: either it has no
proper \(2\)-colouring, or the colours on its nontrivial components
cannot change. Thus \(2\in B_G\). Section 5 shows that \(B_G\) has no
element at least \(n+1\).

By Proposition 4, on every component interval \(I\) of \[
(2,\infty)\setminus\mathcal E_{n+1},
\] either \[
B_G\cap I=\mathbb Q\cap I
\qquad\text{or}\qquad
B_G\cap I=\varnothing.
\] At the finitely many points of \(\mathcal E_{n+1}\), arbitrary
behaviour is possible and harmless.

Consequently \(B_G\) is a union of selected rational open intervals
between consecutive points of \(\mathcal E_{n+1}\), together with a
subset of \(\mathcal E_{n+1}\). It follows that \[
\sup B_G\in\mathcal E_{n+1}.
\]

Finally, \[
\mathfrak M_c(G)=\sup B_G.
\] Every rational \(\ell>\sup B_G\) satisfies the defining universal
mixing condition, while every \(\ell<\sup B_G\) fails it. Whether the
condition holds at equality does not affect the infimum.

This proves the theorem. \(\square\)

\subsubsection{Scope and exceptional
cases}\label{scope-and-exceptional-cases}

The proof does \textbf{not} assume that every orientation class is
connected. That can fail at ratios in \(\mathcal E_n\), where
zero-weight cycles force coordinate differences. Those cases are covered
by retaining all exceptional ratios as possible bad points; their
detailed behaviour cannot create an irrational supremum.

The argument establishes neither that the threshold is always an integer
nor that it is attained.

For edgeless graphs, every colouring reconfiguration graph is connected.
Their threshold is therefore the chosen lower endpoint under the usual
normalization---for example, \(2\) if thresholds are restricted to
\([2,\infty)\). If the displayed definition is read literally with
arbitrarily negative \(\ell\), an edgeless graph instead gives
\(-\infty\); that is a trivial normalization issue in the quoted
statement, separate from the substantive rationality question.

\end{document}
